% operator-stability-v1 — publication template (writeup-aligned)
% Slots filled by scripts/generate_operator_paper.py
\documentclass[12pt]{article}
\usepackage[letterpaper,top=2cm,bottom=2cm,left=3cm,right=3cm]{geometry}
\usepackage[T1]{fontenc}
\usepackage[utf8]{inputenc}
\usepackage{lmodern}
\usepackage{microtype}
\usepackage{amsmath,amssymb,amsthm,mathtools}
\usepackage{booktabs}
\usepackage{xcolor}
\usepackage{enumitem}
\usepackage{xurl}
\usepackage[colorlinks=true,linkcolor=blue,urlcolor=blue,citecolor=blue]{hyperref}

\Urlmuskip=0mu plus 1mu
\urlstyle{tt}
\emergencystretch=2em

\setlength{\parskip}{0.4em}
\setlength{\parindent}{1.2em}
\setlist[itemize]{leftmargin=1.4em,itemsep=0.25em,topsep=0.35em}
\setlist[enumerate]{leftmargin=1.4em,itemsep=0.25em,topsep=0.35em}

\theoremstyle{plain}
\newtheorem{theorem}{Theorem}
\newtheorem{lemma}[theorem]{Lemma}
\theoremstyle{definition}
\newtheorem{definition}[theorem]{Definition}
\newtheorem{example}[theorem]{Example}
\theoremstyle{remark}
\newtheorem{remark}[theorem]{Remark}

% Breakable monospace for Lean paths / identifiers (xurl \path)
\newcommand{\leanpath}[1]{%
  {\raggedright\urlstyle{tt}\path{#1}}}


\definecolor{fundbox}{RGB}{245,245,248}
\definecolor{fundrule}{RGB}{180,180,190}

\title{Median Index Preservation under Bounded Score Perturbations}
\author{Research Architecture Runtime\\[0.25em]
{\normalsize\texttt{operator/median}}}
\date{2026-07-29}

\begin{document}
\maketitle

\begin{abstract}
The median index is the unique strict $k$-th order statistic for the conventional median rank $k$. When every pairwise score gap exceeds $2\varepsilon$, that index is invariant under every additive perturbation with $\|\delta\|_\infty\le\varepsilon$. The claim is Lean-certified ($\texttt{LEAN\_FULL}$) as a definitional reduction of the $k$-th order-statistic margin theorem.
\end{abstract}

\noindent
\begin{center}
\fcolorbox{fundrule}{fundbox}{%
  \begin{minipage}{0.92\linewidth}
  \centering
  \textbf{This theorem is:} \texttt{Reduction}
  \end{minipage}}
\end{center}
\vspace{0.6em}

\section{Problem}
{\raggedright\sloppy
The median operator on a finite score vector $s\in\mathbb{R}^n$ ($n\ge 2$) returns the unique index $i^\star$ that is the strict $k$-th smallest coordinate of $s$, where $k$ is the conventional median rank (so that exactly $k$ scores are strictly smaller than $s_{i^\star}$, and no other index shares the value $s_{i^\star}$). The selected object is this index (equivalently, the unique median position under a strict-order presentation), not a set-valued median band.
\par}

\section{Stability notion}
Admissible perturbations are additive score noise $\delta\in\mathbb{R}^n$ in the closed $\ell_\infty$ ball $\|\delta\|_\infty\le\varepsilon$ for a fixed $\varepsilon\ge 0$. Stability means: if $i^\star$ is the unique strict $k$-th smallest index of $s$ and every pairwise gap of $s$ exceeds $2\varepsilon$, then $i^\star$ remains the unique strict $k$-th smallest index of $s+\delta$ for every such $\delta$. Sharpness: if some rival index lies within distance $2\varepsilon$ of $s_{i^\star}$, an admissible $\delta$ destroys uniqueness of the strict $k$-th position.

\section{Definitions}
\begin{definition}[Strict $k$-th smallest]
Fix $n\ge 2$, $s:\{0,\dots,n-1\}\to\mathbb{R}$, and $k<n$. An index $i$ is the \emph{strict $k$-th smallest} if $\#\{j:s_j<s_i\}=k$ and $s_j=s_i$ forces $j=i$.
\end{definition}

\begin{definition}[All-gaps margin]
Scores satisfy $\mathrm{AllGapsExceed}(s,g)$ if $|s_p-s_q|>g$ whenever $p\neq q$.
\end{definition}

\begin{definition}[$\ell_\infty$ ball]
$\|\delta\|_\infty\le\varepsilon$ means $|\delta_t|\le\varepsilon$ for every coordinate $t$.
\end{definition}

\noindent\textbf{Assumptions.} $n\ge 2$, $\varepsilon\ge 0$, a unique strict $k$-th smallest index exists, and (for invariance) $\mathrm{AllGapsExceed}(s,2\varepsilon)$.

\section{Theorem}
\begin{theorem}[Median / $k$-th margin invariance]\label{thm:inv}
Let $n\ge 2$, $\varepsilon\ge 0$, and let $i$ be the strict $k$-th smallest index of $s$. If $\mathrm{AllGapsExceed}(s,2\varepsilon)$, then for every $\delta$ with $\|\delta\|_\infty\le\varepsilon$, the same index $i$ is the strict $k$-th smallest index of $s+\delta$.
\end{theorem}

\begin{theorem}[Sharpness]\label{thm:sharp}
Under the same standing hypotheses on $(n,s,k,i)$ and $\varepsilon\ge 0$, if there exists $j\neq i$ with $|s_i-s_j|\le 2\varepsilon$, then there exists $\delta$ with $\|\delta\|_\infty\le\varepsilon$ such that $i$ is \emph{not} the strict $k$-th smallest index of $s+\delta$.
\end{theorem}

\section{Intuition}
An $\ell_\infty$ adversary can decrease one coordinate by at most $\varepsilon$ and increase another by at most $\varepsilon$, consuming at most $2\varepsilon$ of any pairwise gap. If every gap is strictly larger than $2\varepsilon$, no pairwise order can reverse and no tie can form, so the count of scores below $s_i$ and the uniqueness of $s_i$ are preserved. If some gap is at most $2\varepsilon$, a midpoint collision forces a tie and uniqueness fails.

\section{Examples}
\begin{example}
Take $n=3$ scores $(s_0,s_1,s_2)=(1,4,7)$ and median rank $k=1$ (unique middle value at index $1$). Pairwise gaps are $3$, $6$, and $3$, so $\mathrm{AllGapsExceed}(s,2\varepsilon)$ holds whenever $2\varepsilon<3$, e.g.\ $\varepsilon=1$. For any $\|\delta\|_\infty\le 1$, the middle index remains $1$: the perturbed values stay strictly ordered. If instead $\varepsilon=2$ so that $2\varepsilon=4$ meets the gap $|4-1|=3\le 4$, the midpoint adversary between indices $0$ and $1$ produces a tie and destroys uniqueness of the strict first order statistic.
\end{example}

\section{Proof}
Write $s'=s+\delta$ and assume $\|\delta\|_\infty\le\varepsilon$.

\paragraph{Pairwise order preservation.}
Fix distinct indices $p,q$ with $|s_p-s_q|>2\varepsilon$. Without loss of generality suppose $s_p<s_q$, so $s_q-s_p>2\varepsilon$. Then
\[
s'_p-s'_q
=
(s_p-s_q)+(\delta_p-\delta_q)
\le
-(s_q-s_p)+2\varepsilon
<
0,
\]
because $|\delta_p-\delta_q|\le 2\varepsilon$. Hence $s'_p<s'_q$. The symmetric argument gives the converse direction: $s'_p<s'_q$ forces $s_p<s_q$. Thus
\[
s_p<s_q\qquad\Longleftrightarrow\qquad s'_p<s'_q.
\]

\paragraph{Count preservation.}
Applying the pairwise equivalence at $q=i$ yields $s_j<s_i\Leftrightarrow s'_j<s'_i$ for every $j$. Therefore
\[
\#\{j:s_j<s_i\}
=
\#\{j:s'_j<s'_i\}.
\]
If $i$ is strict $k$-th smallest for $s$, the left-hand side equals $k$, so the right-hand side equals $k$.

\paragraph{Uniqueness preservation.}
Suppose $s'_j=s'_i$. If $j\neq i$, then $|s_j-s_i|>2\varepsilon$ by the all-gaps hypothesis, and the same $\ell_\infty$ estimate used above yields $s'_j\neq s'_i$, a contradiction. Hence $j=i$.

Combining the two paragraphs shows $i$ remains the strict $k$-th smallest index of $s'$, which is Theorem~\ref{thm:inv}. (Median specialization: take the conventional median rank $k$; the argument never uses any property of $k$ beyond $k<n$.)

\paragraph{Sharpness.}
Let $j\neq i$ satisfy $|s_i-s_j|\le 2\varepsilon$. Define the midpoint collision
\[
\delta_t
=
\begin{cases}
-(s_i-s_j)/2 & t=i,\\
(s_i-s_j)/2 & t=j,\\
0 & \text{otherwise.}
\end{cases}
\]
Then $|\delta_t|\le\varepsilon$ for all $t$ (since $|s_i-s_j|/2\le\varepsilon$), and $s'_i=s'_j$. Uniqueness of a strict $k$-th value at $i$ therefore fails, proving Theorem~\ref{thm:sharp}.

\section{Formal statement}
{\raggedright\sloppy
\begin{description}[style=nextline,leftmargin=1.1em,font=\normalfont\bfseries]
\item[Lean module] \leanpath{Research.Operators.Median.Preservation}
\item[Source file] \leanpath{lean/Research/Operators/Median/Preservation.lean}
\item[Theorems]\leavevmode\par\vspace{0.15em}\begin{itemize}[leftmargin=1.2em]
  \item \leanpath{median\_margin\_invariance}
  \item \leanpath{median\_margin\_sharpness}
\end{itemize}
\item[Note] Definitional reduction of OrderStat.KthMargin.
\item[Certificate] \leanpath{lean/certificates/median/median-margin}
\item[Status] \texttt{LEAN\_FULL} on Mathlib $\mathbb{R}$ (\texttt{REAL\_MATHLIB})
\end{description}
\par}

\section{Proof dependencies}
{\raggedright\sloppy
\begin{itemize}
\item Reduces to \leanpath{Research.Operators.OrderStat.KthMargin} (\leanpath{kth\_margin\_invariance}, \leanpath{kth\_margin\_sharpness}).
\item Uses the $\ell_\infty$ gap estimate $|(\delta_p-\delta_q)|\le 2\varepsilon$ (triangle inequality for absolute values).
\item Uses the definitions of strict $k$-th smallest, all-gaps margin, and midpoint collision adversary from that core.
\item No other operator theorems are required.
\end{itemize}
\par}

\section{Consequences}
{\raggedright\sloppy
Median is the canonical specialization of unique strict order-statistic selection. The same margin theorem governs \texttt{quantile}, \texttt{percentile}, and \leanpath{kth-order-statistic} (all reduce to the same $k$-th core). Ranking and top-$k$ operators use the related pairwise ranking core rather than this count-based $k$-th statement; together they form the order-statistic layer of the primitive library. Any pipeline that selects a median index under bounded score noise may invoke Theorem~\ref{thm:inv} as the stability certificate.
\par}

\end{document}
